% !PRINTED 665
\weberpage{665}
%<<CONT
if $c_1, c_2, \ldots, c_n$ are the numerical values determined by the congruences

\weq{3}{\omega_1 \equiv c_1,\ \omega_2 \equiv c_2,\ \ldots,\ \omega_n \equiv c_n\ (\mathrm{mod}\ \pi)}

By this the values of all integral functions of $z$ at the point $\mathfrak{P}$ are determined. Every integral function $\omega$ is congruent, with respect to the modulus $z - c$, to an expression
\[
a_1\omega_1 + a_2\omega_2 + \cdots + a_n\omega_n
\]
with constant coefficient $a$, and it is then and only then divisible by $\pi$, when
\[
a_1c_1 + a_2c_2 + \cdots + a_nc_n = 0
\]
holds. Accordingly it follows:

2. \gesperrt{Among the integral functions of $z$, only those divisible by $\pi$ have the value 0 at $\mathfrak{P}$.}

Since every fractional function in $\Omega$ can be represented in such a way that numerator and denominator are not both divisible by $\pi$, the value of every function at $\mathfrak{P}$ is thereby determined.

3. \gesperrt{The prime factor $\pi$ is said to be generated by the point $\mathfrak{P}$. It cannot be generated by a point different from $\mathfrak{P}$.}

Conversely, one can also generate a point from every prime factor $\pi$ in $z$.

If $z - c$ is the linear function divisible by $\pi$, then let the function $z$ be given the value $c$. Then take a function $\varrho$ that is divisible by $\pi$, but not by $\pi^2$, and set, in accordance with § 185, (5), for every function $\eta$:

\weq{4}{\eta = a\varrho^m + \eta_1\varrho^{m+1},}

where $m$ is an integral rational number and $\eta_1$ a function whose denominator is not divisible by $\pi$. To the function $\varrho$ let the value 0 be assigned, and to the function $\eta$ the value
\[
\begin{aligned}
\eta_0 &= \infty, \text{ if } m < 0,\\
\eta_0 &= a, \text{ if } m = 0,\\
\eta_0 &= 0, \text{ if } m > 0.
\end{aligned}
\]

4. \gesperrt{The values of all functions $\eta$ so determined satisfy the conditions 1. to 5. and therefore constitute a point, which is said to be generated by the prime functional $\pi$.}

Thus from a single function $z$ and the prime factors of the linear forms $z - c$ one can derive all points at which
%CONT>>
